% =============================================================================
%  مذكرة الأعمال التطبيقية رقم 01 (TP N° 01)
%  المجال 1: بيئة التعامل مع الحاسوب — الوحدة 1: تقنية المعلومات (IT)
%  عنوان الحصة: استكشاف بيئة الحاسوب وتجسيد دورة معالجة البيانات إلى معلومات
%  السنة الدراسية: 2026 / 2027
% =============================================================================
%  البناء (من هذا المجلد):
%    xelatex -interaction=nonstopmode -halt-on-error -output-directory=build tp-note.tex
% =============================================================================

\documentclass[11pt, a4paper]{article}

% ── المشترك ─────────────────────────────────────────────────────────────────
\input{"../../shared/common.tex"}

% ══════════════════════════════════════════════════════════════════════════════
\begin{document}

% ── الصفحة 1: الترويسة + البطاقة الفنية ─────────────────────────────────────
\officialheader

\cardtitle{مذكرة الأعمال التطبيقية رقم: 01 (TP N° 01)}

{\gridsetup
\centering
\begin{tabularx}{0.85\textwidth}{@{}|c|X|@{}}
  \hline
  \cardrow{المؤسسة}{ثانوية الشهيد حكومي العيد أدرار — أدرار}
  \cardrow{الأستاذ}{لفقير عبدالجليل}
  \cardrow{المادة}{المعلوماتية (حصة الأعمال التطبيقية — TP)}
  \cardrow{المستوى والشعبة}{السنة الأولى ثانوي (جذع مشترك علوم وتكنولوجيا / جذع مشترك آداب وفلسفة)}
  \cardrowtwo{المجال التعلمي}{بيئة التعامل مع الحاسوب}{الوحدة التعلمية}{تقنية المعلومات}
  \cardrow{عنوان الحصة التطبيقية}{استكشاف بيئة الحاسوب وتجسيد دورة معالجة البيانات إلى معلومات}
  \cardrow{مكان الإنجاز}{مخبر المعلوماتية}
  \cardrow{المدة}{01 ساعة (لكل فوج)}
  \cardrow{الكفاءة المستهدفة}{القدرة على تشغيل وإيقاف جهاز الحاسوب بسلاسة، استكشاف مواصفاته المادية والنظامية، وتجسيد دورة معالجة البيانات عمليًا.}
  \cardrow{الأهداف التشغيلية}{يطبّق التلميذ خطة التشغيل والإيقاف الآمن لجهاز الحاسوب.\newline
    يستكشف مواصفات العتاد (المعالج، الذاكرة الحية RAM) ونظام التشغيل المباشر من الجهاز.\newline
    يدخل مجموعة بيانات خام ويجري عليها معالجة بسيطة للحصول على معلومات ذات معنى.\newline
    يحفظ العمل في مجلد مخصّص على القرص الصلب.}
  \cardrow{الوسائل المستخدمة}{حواسيب المخبر، نظام التشغيل (Windows)، وبرمجية مكتبية بسيطة (معالج النصوص أو المجدول).}
\end{tabularx}\par
}\vspace{8pt}

% ── الصفحة 2: خطة سير الحصة ────────────────────────────────────────────────
\sectionbreak{خطة سير الحصة التطبيقية}

\tpstage{\badge{1} التمهيد والتوجيهات الأمنية}{10 دقائق}{
  \textbf{التذكير بالدرس النظري:} التذكير السريع بمفهوم البيانات، المعالجة، والمعلومات.\newline
  \textbf{قواعد المخبر والتشغيل:} التوجيه نحو التشغيل الصحيح للوحدة المركزية والشاشة، والتنبيه إلى قواعد السلامة الخاصة بالتوصيلات الكهربائية والاستعمال السليم للفأرة ولوحة المفاتيح.
}

\tpstage{\badge{2} تنفيذ المهام التطبيقية}{40 دقيقة}{
  \textbf{إنجاز المهام الثلاث المبينة في الجدول أسفله، مع متابعة فردية للتلاميذ.}
}

\taskheader
\taskrow{المهمة 1: استكشاف خصائص النظام والعتاد}{
  1. انقر بالزر الأيمن على أيقونة «هذا الكمبيوتر» (Ce PC / This PC) واختر «خصائص» (Propriétés).\newline
  2. حدّد المواصفات التالية وقيدها في ورقة التطبيق:\newline
  \textbullet\ نوع المعالج (Processor).\newline
  \textbullet\ سعة الذاكرة الحية (RAM).\newline
  \textbullet\ اسم ونوع نظام التشغيل المثبّت.
}{
  \textbullet\ ينفّذ التلميذ الخطوات على جهازه.\newline
  \textbullet\ يستخرج مواصفات الجهاز وقدراته التخزينية والمعالجية.
}

\taskrow{المهمة 2: تجسيد دورة معالجة البيانات إلى معلومات}{
  1. افتح محررًا نصيًا أو تطبيقًا مكتبيًا بسيطًا.\newline
  2. \textbf{إدخال البيانات الخام (Data):} قم بكتابة البيانات التالية غير المرتبة:\newline
  \phantom{0}[علي — 2008 — 14.5]، [سارة — 2009 — 09.0]، [عمر — 2008 — 11.0].\newline
  3. \textbf{إجراء المعالجة (Processing):}\newline
  \textbullet\ قم بترتيب الأسماء أبجديًا.\newline
  \textbullet\ فرز التلاميذ حسب الملاحظة (ناجح: معدل $\ge 10$ / راسب: معدل $< 10$).\newline
  4. \textbf{استخراج المعلومة (Information):} صياغة التقرير النهائي (مثال: قائمة الناجحين ومعدل الفوج).
}{
  \textbullet\ يدخل التلميذ البيانات الخام عبر لوحة المفاتيح.\newline
  \textbullet\ يطبّق عملية الترتيب والفرز المنطقي.\newline
  \textbullet\ يستنتج كيف تحوّلت الأرقام المجرّدة إلى معلومات مفيدة.
}

\taskrow{المهمة 3: تنظيم الملفات والتخزين}{
  1. أنشئ مجلدًا جديدًا على القرص (D:) باسم الفوج والتلميذ.\newline
  2. احفظ الملف داخل المجلد باسم TP01\_Informatique.\newline
  3. أعد إغلاق البرنامج والتأكد من وجود الملف في مساره.
}{
  \textbullet\ يتدرب على إنشاء المجلدات وحفظ الملفات في وسائط التخزين لحمايتها.
}

\vspace{6pt}
\tpstage{\badge{3} التقييم والإغلاق الآمن}{10 دقائق}{
  \textbf{المراقبة والمصادقة:} يمرّ الأستاذ على الحواسيب للتحقق من إنجاز المهام الثلاث وحفظ الملفات في مسارها الصحيح.\newline
  \textbf{مناقشة صعوبات التنفيذ:} معالجة الأخطاء الشائعة (مثل عدم تحديد مسار الحفظ، أو الخطأ في الوصول لخصائص النظام).\newline
  \textbf{الإغلاق الآمن:} توجيه التلاميذ لإغلاق الحاسوب عبر أمر إيقاف التشغيل (Arrêter / Shut down) ثم ترتيب مكان العمل.
}

\end{document}
